\scanpage{15}
\uline{Exercice}\qquad Réalisation régulière d'un polygône.

Soit $\deleted{V}E$ plan affine euclidien sur $\mathbb{R}$ (avec \unsure{metrique|métrique} déf pos. sur \unsure{l'espace|l'esp|l'esp.} $V$ des translations). Une réalisation rectiligne d'un polygone $(S,A,R)$ est dite \uline{régulière} si tous les cotes du polygône ont $=$
\tikz[baseline=(longueur.base)]{
  \node[anchor=base west,inner sep=0pt] (longueur) at (0,0) {longueur};
  \node[align=left,inner sep=0pt] at (3.25,.20) {\inserted{\shortstack[l]{et tous les angles\\itou}}};
  \draw (1.70,.75) .. controls (2.1,.82) and (4.6,.74) .. (4.90,.69)
    .. controls (5.17,.12) and (4.86,-.38) .. (4.15,-.49)
    .. controls (3.70,-.62) and (3.20,-.24) .. (2.98,-.22)
    .. controls (2.45,-.17) and (2.12,.27) .. (1.83,-.03)
    -- (1.59,-.29)
    .. controls (1.90,.15) and (1.90,.28) .. (1.65,.51)
    .. controls (1.40,.72) and (1.62,.74) .. (1.70,.75);
}, régulière \deleted{et} \uline{\unsure{normalisee|normalisée}} si la dite longueur est 1.

a) Existence, grâce à équation cyclotomique $z^n=1$ dans $\mathbb{C}$.

b) Prouver que si $(V\deleted{\unsure{,}}\supset P\deleted{\unsure{,}}\supset\deleted{\unsure{A}}S)$ et $(V'\supset P'\supset S')$ sont deux $n$-polygônes ($n\geq3$) munis d'une réal. rectiligne \uline{rég.} (resp \unsure{normalisee|normalisée}) il existe une \deleted{\unsure{iso}} similitude (resp \unsure{isom}) \deleted{\uline{unique}} de $V$ sur $V'$, qui transforme $P$ et $P'$, $S$ en $S'$; \deleted{si $V\unsure{;|,}P,S$} plus précisément, pour tt isom (comb) de $(\deleted{\unsure{P}}S,A,R)$ sur $(S',A',R')$, $\exists\,\unsure{\vert|1}$ similitude
% Raised insertion with its curved, right-pointing connector above de V.
\tikz[baseline=(deV.base)]{
  \node[anchor=base west,inner sep=0pt] (deV) at (0,0) {de $\deleted{\unsure{E}}V$};
  \node[anchor=south west,inner sep=0pt] at (-.42,.40) {\inserted{(isom)}};
  \draw[->] (-.18,-.05) .. controls (.05,.42) and (1.10,.54) .. (1.55,.57);
}
sur $V'$ qui l'induise. \deleted{Montrer que}

\deleted{\unsure{Ex|Cor}}\quad\uline{Cor} Notion d'enveloppe affine (ou vectorielle, grâce : centre de
% The insertion stands above gravité), between gravité and the closing parenthesis.
\tikz[baseline=(gravite.base)]{
  \node[anchor=base west,inner sep=0pt] (gravite) at (0,0) {gravité)};
  \node[anchor=south west,inner sep=0pt] at (-.10,.50) {\inserted{de l'ens des sommets}};
  \draw (1.44,-.26) .. controls (1.58,.26) and (1.77,.46) .. (2.08,.47)
    .. controls (2.76,.49) and (3.60,.39) .. (3.86,.48)
    -- (4.07,.66);
  \draw (3.86,.48) -- (3.91,.70);
}
d'un polygône combinatoire.

\deleted{\unsure{Si|Soit}} Orientation du polygône $\longleftrightarrow$ orientation de son enveloppe

% Two detached dark points below the right-hand end of the last line.
% Writing versus surface traces is unresolved; no punctuation value is assigned.
\figure{\begin{tikzpicture}[x=1cm,y=1cm]
  \path[use as bounding box] (0,0) rectangle (13.5,1.25);
  \node[inner sep=0pt] at (12.4,.20) {\unsure{\tikz[baseline=-.3ex]{
    \fill (.055,.08) ellipse [x radius=.018cm,y radius=.024cm];
    \fill (0,0) ellipse [x radius=.026cm,y radius=.016cm];
  }}};
\end{tikzpicture}}
% archive sheet number: 14
